\chapter{Theory}

Theory
Studies of charity advertising have primarily focused on “poverty porn,” images intended to shock audiences into action by showing extreme sickness, hunger, poverty, etc. The name “poverty porn” was popularized during the 1980s as a critique of the widespread imagery of the Ethiopian famine (showcasing severe illness and emaciation) as being exploitative and voyeuristic. The images received criticism for their demeaning portrayal of beneficiaries and reinforcement of harmful stereotypes, including colonialist narratives about the backwardness of the Global South, leading the General Assembly of European NGOs to adopt a code of conduct in 1989 that discouraged NGOs from using “pathetic images” (Pedersen, 2014). Despite these guidelines, the charity sector persisted in using sensationalized imagery. 

Charities justified the continued prevalence of “poverty porn” by arguing that less sensationalized advertising may reduce the amount that charities fundraise, limiting their capabilities. For example, images of children intended to evoke positive emotions resulted in fewer potential donations than those intended to evoke negative emotions (Burt \& Strongman, 2005). Similarly, sad expressions garner more donations than happy or neutral expressions (Small \& Verrochi, 2009). Beyond the emotion itself, its intensity is also relevant; advertisements with higher emotional intensity lead to more donations (Pansoni \& Gistri, 2025). These studies suggest the efficacy of shock-based advertising, which typically relies on sad images with high emotional intensity. The success of this advertising strategy has been corroborated by some International Non-Governmental Organizations (INGOs), which raised millions of dollars through these campaigns; for example, UK donors reported to Save the Children that they were more willing to give money after being shown shocking images (Burt \& Strongman, 2005).

Despite the immediate benefits of sensationalized advertising, it has been criticized for prioritizing short-term goals at the expense of long-term systemic change (Tallon, 2008). Overexposure to shock-based advertising can result in “donor fatigue,” leading people in the Global North to feel that the situation in the Global South is beyond improvement, or to disengage from the misery and guilt these advertisements evoke (Chouliaraki, 2010). Furthermore, if the audience believes the campaign's representation of the Global South is misrepresentative, it can lead to INGOs receiving less financial support (Girling \& Adesina, 2024; Ademolu \& Crombie, 2024). Additionally, a study of residents of low-income settlements in Kenya found that the narrative most likely to promote community skill-building did not decrease the likelihood of donation (Thomas et al., 2020). The same study found that beneficiaries' outcomes are influenced by more than just the amount of money received; participants had higher levels of self-efficacy and anticipated social mobility when receiving money from a charity with an empowerment narrative rather than a poverty-alleviation narrative (Thomas et al., 2020). The notion that the amount of money raised does not necessarily correspond to the best outcome for beneficiaries is a concern with shock-based advertising, which, by poorly representing regions, may reduce the investment they receive and disincentivize tourism (Tallon, 2008).

Negative representations of the Global South in charity advertising also influence public perception in the Global North, a phenomenon that has been most researched in regards to Africa. Charity advertising contributes to the widespread belief that Africa is composed of primitive tools, homelessness, hunger, poverty, and wild animals (Ademolu, 2021). A 2012 Oxfam survey found that 55\% of respondents' first association with Africa was a problem such as famine, hunger, or poverty (Pedersen, 2014). The conception of Africa as backward or inferior also leads to the idea that it is reliant on the Global North. 71\% of UK adults believed that the Global South depended on the money and knowledge of the West to progress according to a 2001 Volunteer Work Abroad (VSO) survey (Pedersen, 2014). Negative representations of Africa have been found to lead respondents to perceive those in the developing world as having less agency, and recipients perceived as having less agency are consulted less when considering how funds are used (Clough et al., 2023). Furthermore, the internalization of negative ideas about Africa also extends to diaspora communities. Charity advertisements made some members of black diaspora communities in Britain more reluctant to claim their African identity because they felt that the advertisements portrayed Africans and blackness as being inferior (Ademolu, 2021). 

The Global North’s portrayal of Global South members as subordinate to and dependent on the Global North is a key component of postcolonial theory, in particular Edward Said’s Orientalism, which suggests that the colonizer defines itself in relation to the colonized. Within this theoretical framework, the West conceptualizes its own superiority by comparing itself to the inferior “Other”. Colonial binaries define this opposition: if the West views itself as “white”, “masculine”, and “civilized”, then it views the “Other” as “black”, “feminine”, and “barbaric” (Ademolu \& Crombie, 2024; Konaté, 2020). Furthermore, the West justifies its domination of the “Other” by portraying it as impoverished, backward, and helpless (Matthew, 2021). This colonial narrative has been applied to Africa, in discourse that suggests it is a “dark continent" characterized by war, barbarism, and suffering (Oguh, 2015). This treats all of Africa as though it is a homogenous region composed of tribal hunter-gatherers and “plagued by meaningless violence”  (Oguh, 2015; Tallon, 2008; Bonsu, 2009). The narrative of violence that is meaningless and poverty that is merely a lack of material resources masks their systematic causes, including illiteracy, war, and corruption (Al-Sabbagh, 2023). Framing violence in Africa as senseless is used as justification for another aspect of the colonial narrative: that the colonized need the colonizer to bring modernity and civilization (Tallon, 2008; Bonsu, 2009). Historically, this narrative was religiously championed as it was the “white man’s burden” to save the colonized from themselves and bring them towards salvation (Tallon, 2008; Bonsu, 2009). The idea of the Global North as supporting a Global South that depends on its aid has been reinforced by charity advertisements.

Charity advertisements have been criticized for contributing to colonialist discourse by dismissing examples of modernity and infrastructure in the Global South. For example, charity advertisements frequently showcase rural areas and traditional dress (Bonsu, 2009). This excludes symbols of development present in indigenous media depictions, including local health facilities and social services, which charities are theoretically working in tandem with. The invisibility of local efforts at community and societal improvement makes the Global South seem like a passive party that depends on the charity of the Global North. In charity advertising, members of the Global North are commonly portrayed as the active parties: doctors, home builders, etc. In contrast, members of the Global South are more likely to be portrayed as the passive party, looking desolate and waiting for help to arrive (Pedersen, 2014). 

The Global North/Global South binary affects not only how issues are framed but also which issues are represented. This divide was highlighted by a 2012 campaign targeting white British poverty, entitled “It Shouldn’t Happen Here,” which was criticized by centrist and right-wing commentators for the notion that people in the Global North could be experiencing poverty equivalent to that experienced in the Global South and that the comparison was politicizing poverty (Kelbert, 2022). This reflects the idea that extreme poverty in the Global North either does not exist or constitutes a political or personal failure, which influences not only how issues are framed but also the nature of the issues covered. Alternatively, those experiencing poverty in the Global South are seen as victims of cultures that perpetuate poverty, making them “deserving” of charity (Kelbert, 2022). Advertisements targeting the Global South tend to emphasize basic survival needs, while those targeting the Global North tend to emphasize specialized institutional needs or medical services.

This victimization of the Global South has been emphasized through the portrayal of “ideal victims”, particularly women and children, who are seen as vulnerable, respectable and likely to command societal sympathy. A widely distributed image in charity advertising was that of the ‘de-contextualized lone child’ (Tallon, 2008). When images of a single child alone and suffering are removed from explanations about its cause, all of Africa becomes represented by the suffering child in the eyes of the West (Tallon, 2008). This use of children in charity advertising has been linked to the infantilization of the other and the paternalization of the West (Pedersen, 2014). Besides children, another “ideal victim” to inspire donations is women, who are often portrayed occupying traditional roles and accompanied by their children, while men are often missing and, when portrayed, are often shown as corrupt leaders or guerrilla fighters (Dogra, 2011). This combined narrative is that women in the Global South are unlike modern women in the Global North and Western ideals need to be imposed onto them to save them from the oppression of their male counterparts (Martínez, 2024).  

While charity advertising representing the Global South as suffering and in need of aid from the Global North has been an established trend, it is unclear how this has changed as it shifts from a humanitarian framework to a post-humanitarian one. Humanitarianism is the advertising strategy that has historically defined charity advertising; this method persuades donors to give by appealing to their emotions, either negative emotions (e.g. sadness, guilt, indignation) or positive emotions (e.g. gratitude) (Chouliaraki, 2010). Both negative and positive humanitarianism have received criticisms, negative humanitarianism for victimizing beneficiaries and positive humanitarianism for masking real-life imbalances in favor of a happier narrative (Turton, 2013; Chouliaraki, 2010). To address the pitfalls of humanitarianism, there has been a transition to post-humanitarian advertising, which appeals to the logic of benefits achievable through donations, making charity appeals less political and more commercial (Gill, 2020; Martínez, 2024). However, post-humanitarianism may be reshaping rather than eliminating colonial narratives, focusing on the role of the “self” in ways that perpetuate the charity-and-donor-as-saviors narrative (Martínez, 2024). Post-humanitarianism has been critiqued as driven by neoliberal ideas, in which donors can buy ethical goodness without substantive thought about the issues (Chouliaraki, 2010). This prompts the need for further research into the impact of changing advertising trends on colonial narratives.

Simultaneously, charity advertisements have undergone another transformation, with improved representation of the Global South in imagery. Following the Black Lives Matter movement, many charities updated their image guidelines, including Oxfam's Content Guidelines (2020), WaterAid's Ethical Image Policy (2021), and USAID's Photo Guide for Partners (2020) (Girling \& Adesina, 2024). There have also been more discussions about power shifting, allowing the marginalized communities represented to have sway over the decisions made and the narratives told (Ademolu \& Crombie, 2024). This was especially relevant when the COVID-19 pandemic led to the hiring of local photographers (Girling \& Adesina, 2024). These changes led to fewer pitiful images being used in UK newspaper advertisements; however, there remains an underrepresentation of urban settings, an overrepresentation of women and children as victims, and a tendency to portray beneficiaries as passive (Girling \& Adesina, 2024). It is unclear whether the improvements in representation extended to the advertisements' text.

While the text of charity advertisements has been understudied, evidence suggests it may serve an important function. In this semiotic analysis framework, the images and text both serve as signs that together constitute the advertisement's overall message, which includes the general meaning (denotation) and cultural or emotional meaning (connotation) (Sudi, 2024). Furthermore, the text accompanying the image may serve as what Roland Barthes termed “anchorage”, a determinant of the meaning the audience ascribes to the image (Barthes, 1977). In the context of charity advertising, a single image can take on entirely different connotations depending on the advertising text, potentially undermining improvements in imagery. This is especially critical because charities have an incentive to use pitiful imagery, as it has proved a successful advertising strategy. Still, that practice has been mitigated due to updated image guidelines and public outcry against pitiful imagery. However, the text has the power to recontextualize images that are not inherently pitiful, and it has not been subjected to the same scrutiny as the images.

Furthermore, the text of advertisements is critical to analysis because people have been shown to unconsciously adopt the language of dominance (Bonsu, 2009). This aligns with Foucault’s theories about representation and discourse, namely that discourse governs how topics are thought about and that the knowledge and power produced by these discourses can have real systemic impacts (Hall 1997). In this context, with a discourse surrounding the Global South as impoverished, suffering, and helpless, it can impact the way the Global South is treated, contributing to systemic inequality that reinforces the original discourse.

Despite the role text can play in shaping the interpretation of images and the messages of charity advertisements, it has received less public censure and thus less oversight. There have been some updated guidelines for language; notably, Oxfam’s Inclusive Language Guide specifically references "eliminat[ing] colonial legacies,” but it is unclear to what extent these measures have been adopted in practice (Oxfam, 2023). One study noted that while images portray subjects more actively, the captions do not follow suit, but no quantitative evidence exists to substantiate this claim (Girling \& Adesina, 2024). Due to the historical propensity toward colonial language and the limited outrage directed at the text compared to the images, this study hypothesizes that the text of charity advertisements is likely to continue depicting victims as vulnerable despite improvements in images. Vulnerability language is not exclusive to Global South advertisements; regardless of geographic location, charities have an incentive to make the need for donations seem imperative by portraying recipients as helpless. However, if vulnerability language is motivated by colonial narratives, there will be a higher density of vulnerability themes in advertisements featuring Global South beneficiaries than in those featuring Global North beneficiaries.

H1. Advertisements featuring Global South beneficiaries will use a higher proportion of words related to vulnerability than those featuring Global North beneficiaries. 

Existing research suggests that the shift from humanitarianism to post-humanitarianism is not a unilateral move away from colonial narratives in writing, but rather a change in how these narratives are represented (Gill, 2020). The language in charity advertisements has been criticized as “paternalistic” and “othering,” with these ideas often conveyed subtly (Coffey, 2023). An example of this was a comparison between the phrases “lift someone out of poverty” and “serve communities seeking to eliminate poverty” (Coffey, 2023). The former phrase portrays the charity as the actor taking the recipient out of poverty, while the beneficiary is a passive recipient of aid. The latter phrase portrays the community as another active party that is seeking to eliminate poverty. This language critique aligns with posthumanitarian theories that, in the absence of emotional advertising, there will be a greater focus on the self. This also agrees with post-colonial theories suggesting that the Global North parties (such as the charities and donors) will be portrayed as active in contrast to the inactivity of Global South actors (the beneficiaries). As with the vulnerability analysis, charities being shown with agency and beneficiaries being shown without are not inherently indicative of colonial narratives, and are likely to be present across regions to highlight charity accomplishments and beneficiary needs; to determine if they are being motivated by colonialism, it is necessary to compare the proportion of language for Global South advertisements to the proportion for Global North advertisements.

H2: Compared to advertisements with Global North beneficiaries, advertisements with Global South beneficiaries will have a higher proportion of words associated with charity and donor activity.

This study anticipates that, given historical differences in the representation of the Global South, a charity advertisement focusing on beneficiaries in the Global South is more likely to include language associated with colonial narratives. This applies to overt victimizing language that emphasizes the beneficiary's helplessness, which has been criticized in negative humanitarian advertisements. This study also expands the analysis to consider the amount of agency language associated with the charity and the donor, as post-humanitarian advertisements have been criticized for giving more agency to the charity/donor. To test these theories, this study performed a cross-sectional quantitative comparison of the normalized thematic word counts of theoretically significant language in charity advertisements.
